% !PRINTED 540
\weberpage{540}

%<<CONT
wenn $\omega$ ein vollständiges Wurzelsystem von Formen mit ungeraden
äußeren Koeffizienten durchläuft, mit Benutzung der Formel:
\[
f_1(\omega')f_2(\omega')f(\omega') = \sqrt{2},
\]
% [sic] The source prints -m' here but defines h' for -m immediately below (23).
wenn $h'$ die zur Diskriminante $-m'$ gehörige Klassenzahl ist:
% [sic] The middle numerator in (23) is printed 2^h; the product has h' pairs, so consistency with the last equality requires 2^{h'}.
\weq{23}{\Pi f(\omega) = \Pi f_1(\omega_1)f_2(\omega_2) = \frac{2^h}{\Pi f_1(\omega')f_2(\omega')} = \sqrt{2}^{\,h'} \Pi f(\omega'),}
worin $h$, $h'$ die Klassenzahlen für die Diskriminanten $-4m$, $-m$
bedeuten. Es ist dann (§ 123):
\[
h = 2h'.
\]

Sind also wie oben $\tau$, $\tau'$ so bestimmt, daß
\[
\frac{f(\omega)}{2^\tau}, \quad \frac{f(\omega')}{2^{\tau'}}
\]
Einheiten werden, so ist
\[
\Pi f(\omega) = 2^{h\tau}, \quad \Pi f(\omega') = 2^{h'\tau'},
\]
und daraus nach (23):
\weq{24}{\tau = \frac{\tau'}{2} + \frac{1}{4},}
eine Formel, die auch noch für $m' = 1$ gilt, wo $h' = h$ und
die beiden Werte $2\,\omega'$ und $\omega' : 2$ äquivalent sind.

Durch wiederholte Anwendung dieser Formel ergibt sich, wenn
\[
m = 4^\lambda m'
\]
ist, für ein beliebiges positives $\lambda$:
\[
\tau = \frac{\tau'}{2^\lambda} + \frac{1}{2} - \frac{1}{2^{\lambda+1}}.
\]

Fassen wir das hiermit Bewiesene zusammen, so können wir
sagen, daß folgende Größen \gesperrt{algebraische Einheiten} sind:
\[
\begin{aligned}
\frac{f(\omega)}{\sqrt[4]{2}}, & \quad m \equiv 1, 2 \pmod 4, \\
\frac{f(\omega)}{\sqrt[3]{2}}, & \quad m \equiv 3 \pmod 8, \\
\frac{f(\omega)}{\sqrt{2}}, & \quad m \equiv 7 \pmod 8,
\end{aligned}
\]
%CONT>>
